\documentclass[UTF8,11pt]{ctexart}
\usepackage[a4paper,margin=22mm]{geometry}
\usepackage{graphicx,enumitem,amssymb,tabularx}
\usepackage[normalem]{ulem}
\usepackage{xeCJKfntef}
\setlength{\parindent}{0pt}
\setlength{\parskip}{.6em}
\sloppy
\begin{document}
\section*{2023年12月英语六级真题(第2套)}
\subsection*{2023 年12 月大学英语六级考试真题（第2 套）}


\subsection*{\textbf{Part I Writing} \textbf{(30 minutes)}}


\textbf{Directions:}\textit{ For this part, you are allowed 30 minutes to write an essay that begins with the sentence }\textbf{\textit{“Nowadays}} \textbf{\textit{parents are increasingly aware that allowing kids more freedom to explore and learn on their own helps foster}} \textbf{\textit{their independence and boost their confidence.”}}\textit{ You can make comments, cite examples or use your personal} \textit{experiences to develop your essay. You should write at least }\textit{\uline{150 }}\textit{words but no more than }\textit{\uline{200 }}\textit{words (}\textit{\uline{not including}} \textit{\uline{the sentence given}}\textit{).}


\subsection*{\textbf{Part II} \textbf{Listening Comprehension} \textbf{(30 minutes)}}


\subsection*{\textbf{Section A}}


\textbf{Directions:} \textit{In this section, you will hear two long conversations. At the end of each conversation, you will hear} \textit{four questions. Both the conversation and the questions will be spoken only once. After you hear a question, you} \textit{must choose the best answer from the four choices marked A), B), C) and D). Then mark the corresponding letter} \textit{on }\textbf{\textit{Answer Sheet 1}}\textit{ with a single line through the centre.}


\textbf{Questions 1 to 4 are based on the conversation you have just heard.}


1.


\begin{itemize}[leftmargin=2em,itemsep=.2em,topsep=.4em]\item[A.] Spending their holidays in a novel way.
\item[B.] Exploring more summer holiday resorts.
\item[C.] Surfing online to check out the best deals.
\item[D.] Renting a car instead of driving their own.
\end{itemize}

2.


\begin{itemize}[leftmargin=2em,itemsep=.2em,topsep=.4em]\item[A.] He did not like to be locked into one place.
\item[B.] He did not find holiday homes appealing.
\item[C.] He was well travelled.
\item[D.] He once owned a van.
\end{itemize}

3.


\begin{itemize}[leftmargin=2em,itemsep=.2em,topsep=.4em]\item[A.] Generate their own electricity.
\item[B.] Drive under any weather condition.
\item[C.] Receive instructions via computers.
\item[D.] Ensure the safety of passengers.
\end{itemize}

4.


\begin{itemize}[leftmargin=2em,itemsep=.2em,topsep=.4em]\item[A.] Riding one’s mountain bike on vacation.
\item[B.] Slowing down in one’s increasingly hectic life
\item[C.] Enjoying the freedom to choose where to go and work.
\item[D.] Having one’s basic needs covered while away from home.
\end{itemize}

\textbf{Questions 5 to 8 are based on the conversation you have just heard.}


5.


\begin{itemize}[leftmargin=2em,itemsep=.2em,topsep=.4em]\item[A.] Her physical health has deteriorated these past few weeks.
\item[B.] She has lagged behind most of her co-workers in output.
\item[C.] Her job performance has worsened over the past month.
\item[D.] She has missed several important appointments lately.
\end{itemize}

6.


\begin{itemize}[leftmargin=2em,itemsep=.2em,topsep=.4em]\item[A.] Penalty for curtailed output.
\item[B.] Some problems at home.
\item[C.] Disturbance of her mind.
\item[D.] Serious health issues.
\end{itemize}

7.


\begin{itemize}[leftmargin=2em,itemsep=.2em,topsep=.4em]\item[A.] The woman’s whole- hearted support.
\item[B.] The woman’s work proficiency.
\item[C.] His management capability.
\item[D.] His engaging personality.
\end{itemize}

8.


\begin{itemize}[leftmargin=2em,itemsep=.2em,topsep=.4em]\item[A.] The man will help the woman get back to her usual self.
\item[B.] The man will be back at his 100\% in a couple of weeks.
\item[C.] The woman will be off work on the next two Mondays.
\item[D.] The woman will resume her work in two weeks.
\end{itemize}

\subsection*{\textbf{Section B}}


\textbf{Directions:} \textit{In this section, you will hear two passages. At the end of each passage, you will hear three or four} \textit{questions. Both the passage and the questions will be spoken only once. After you hear a question, you must} \textit{choose the best answer from the four choices marked A), B), C and D). Then mark the corresponding letter on} \textbf{\textit{Answer Sheet 1}}\textit{ with a single line through the centre.}


\textbf{Questions 9 to 11 are based on the passage you have just heard.}


9.


\begin{itemize}[leftmargin=2em,itemsep=.2em,topsep=.4em]\item[A.] It can have an impact on our moods and emotions.
\item[B.] It can improve our financial status significantly.
\item[C.] It can help us achieve better work performance.
\item[D.] It can enable us to live a healthier and longer life.
\end{itemize}

10.


\begin{itemize}[leftmargin=2em,itemsep=.2em,topsep=.4em]\item[A.] One’s health tends to differ before and after marriage.
\item[B.] The spouse’s level of education can impact one’s health.
\item[C.] The wealthier one’s spouse is, the healthier one becomes.
\item[D.] One’s health status is related to one’s social background.
\end{itemize}

11.


\begin{itemize}[leftmargin=2em,itemsep=.2em,topsep=.4em]\item[A.] They had more education than their spouses.
\item[B.] They had much in common with their spouses.
\item[C.] They benefited a lot from their career achievements.
\item[D.] They showed interest in their spouses’ occupations.
\end{itemize}

\textbf{Questions 12 to 15 are based on the passage you have just heard.}


12.


\begin{itemize}[leftmargin=2em,itemsep=.2em,topsep=.4em]\item[A.] Finding out the changing climate patterns.
\item[B.] Identifying the cities’ geological features.
\item[C.] Forecasting flood risks accurately.
\item[D.] Eliminating their root cause.
\end{itemize}

13.


\begin{itemize}[leftmargin=2em,itemsep=.2em,topsep=.4em]\item[A.] To validate his hypothesis about the gravity of floods.
\item[B.] To determine the frequency of high tides causing floods.
\item[C.] To see the feasibility of his project on flooding.
\item[D.] To improve his mathematical flooding model.
\end{itemize}

14.


\begin{itemize}[leftmargin=2em,itemsep=.2em,topsep=.4em]\item[A.] To forecast rapid floods in real time.
\item[B.] To classify the flooding data processed.
\item[C.] To study the consequences of high tides on flooded areas.
\item[D.] To teach local citizens how to collect data of incoming floods.
\end{itemize}

15.


\begin{itemize}[leftmargin=2em,itemsep=.2em,topsep=.4em]\item[A.] They tracked the rising tides with video-cameras.
\item[B.] They set up Internet-connected water-level sensors.
\item[C.] They used newly-developed supercomputing facilities.
\item[D.] They observed the direction of water flow on the spot.
\end{itemize}

\subsection*{\textbf{Section C}}


\textbf{Directions:}\textit{ In this section, you will hear three recordings of lectures or talks followed by three or four questions.} \textit{The recordings will be played only once. After you hear a question, you must choose the best answer from the four} \textit{choices marked A), B), C) and D). Then mark the corresponding letter on }\textbf{\textit{Answer Sheet 1}}\textit{ with a single line} \textit{through the centre.}


\textbf{Questions 16 to 18 are based on the recording you have just heard.}


16.


\begin{itemize}[leftmargin=2em,itemsep=.2em,topsep=.4em]\item[A.] To debate the validity of current survey data.
\item[B.] To argue about the value of a college degree.
\item[C.] To account for the drastic decline in employment among men.
\end{itemize}

\begin{itemize}[leftmargin=2em,itemsep=.2em,topsep=.4em]\item[D.] To compare men without college degrees with those who have.
\end{itemize}

17.


\begin{itemize}[leftmargin=2em,itemsep=.2em,topsep=.4em]\item[A.] The issue of changing job requirements.
\item[B.] The increase in women taking up jobs.
\item[C.] The impact of inflation.
\item[D.] The factor of wages.
\end{itemize}

18.


\begin{itemize}[leftmargin=2em,itemsep=.2em,topsep=.4em]\item[A.] The sharp decline in marriage among men with no college degrees.
\item[B.] The wage gap between those with college degrees and those without.
\item[C.] More jobs requiring their holders to have a college degree nowadays.
\item[D.] Men’s unwillingness to accept low wages in times of growing inflation.
\end{itemize}

\textbf{Questions 19 to 21 are based on the recording you have just heard.}


19.


\begin{itemize}[leftmargin=2em,itemsep=.2em,topsep=.4em]\item[A.] More and more people attach importance to protecting endangered animals.
\item[B.] An increasing number of people demand to free animals being kept in cages.
\item[C.] More and more people prioritize animal welfare when buying things to wear.
\item[D.] An increasing number of people follow the latest trend of becoming vegetarians
\end{itemize}

20.


\begin{itemize}[leftmargin=2em,itemsep=.2em,topsep=.4em]\item[A.] Utilized a silk substitute made from mushrooms.
\item[B.] Refrained from using chemicals in their products.
\item[C.] Labelled all their products as vegan.
\item[D.] Avoided the use of leather and fur.
\end{itemize}

21.


\begin{itemize}[leftmargin=2em,itemsep=.2em,topsep=.4em]\item[A.] Whether they can be regarded as ethical.
\item[B.] Whether they can be considered sustainable.
\item[C.] Whether they actually signify a substantial change.
\item[D.] Whether they effectively protect animals at large.
\end{itemize}

\textbf{Questions 22 to 25 are based on the recording you have just heard.}


22.


\begin{itemize}[leftmargin=2em,itemsep=.2em,topsep=.4em]\item[A.] The belief is less prevalent that the world is going to hell.
\item[B.] The environmental welfare has worsened in the world.
\item[C.] The world has seen more violence in recent years.
\item[D.] The era we live in is the most peaceful in history.
\end{itemize}

23.


\begin{itemize}[leftmargin=2em,itemsep=.2em,topsep=.4em]\item[A.] They did not wish to live in the previous century.
\item[B.] They were convinced by the statistics presented to them.
\item[C.] They believed the world was deteriorating.
\item[D.] They were actually not in their right mind.
\end{itemize}

24.


\begin{itemize}[leftmargin=2em,itemsep=.2em,topsep=.4em]\item[A.] Our ancestors’ influence.
\item[B.] Our psychological biases.
\item[C.] The current state of affairs.
\item[D.] The subjectivity of mass media.
\end{itemize}

25.


\begin{itemize}[leftmargin=2em,itemsep=.2em,topsep=.4em]\item[A.] Paying attention to negative information.
\item[B.] Calculating dangerous risks to our survival.
\item[C.] Vacuuming up depressing or enraging stories.
\item[D.] Spreading exciting news around us far and wide.
\end{itemize}

\subsection*{\textbf{Part III Reading Comprehension (40 minutes)}}


\subsection*{\textbf{Section A}}


\textbf{Directions:} \textit{In this section, there is a passage with ten blanks. You are required to select one word for each blank} \textit{from a list of choices given in a word bank following the passage. Read the passage through carefully before} \textit{making your choices. Each choice in the bank is identified by a letter. Please mark the corresponding letter for} \textit{each item on }\textbf{\textit{Answer Sheet 2}}\textit{ with a single line through the centre. You may not use any of the words in the bank} \textit{more than once.}


The desert is deceiving. At first glance it looks lifeless, barren, and bone-dry. For most passersby humming through the Mojave on their way to try their luck in Las Vegas or heading towards the Grand Canyon, it’s just a \uline{26} stretch of land with some mountains in the distance and more \uline{27} to be a setting for a movie that takes place on Mars. The desert, however, is \uline{28} with life, mystery, and magic.


The Mojave desert sees less than two inches of rain a year, and like most deserts, is a land of \uline{29 }. Temperature fluctuations vary from freezing to \uline{30} hot, not only between seasons, but even within the same day.


At the heart of this vast, 25,000-square-mile desert is Mojave National Preserve. The folks managing and working there wear the wide-brimmed (宽边的) hats that have become \uline{31} with places like Yosemite. The preserve is rich with history, culture, biodiversity, and endless opportunities to \uline{32} your interest, especially for anyone who loves the outdoors. For a photographer, it requires patience and \uline{33} a few spare tires to maneuver the network of unpaved “roads”, but the rewards are plenty.


Photographing here requires a good, \uline{34} pair of shoes and a lot of preparation. Plenty of water, spare tires, and a full tank of gas are needed to explore this unit of the national park system, which is the third largest in the country.


Like all deserts. though, if you are patient, you will be rewarded, as they often reveal their secrets slowly. During my last outing, I was surprised to see how much of the area was shaped by ancient volcanoes and geological forces, much of which remains today, giving the area a \uline{35} feel and painting a colorful background for great photographs.


\begin{itemize}[leftmargin=2em,itemsep=.2em,topsep=.4em]\item[A.] apt
\item[B.] burning
\item[C.] extremes
\item[D.] flat
\item[E.] fractions
\item[F.] overflowing
\item[G.] parasites
\item[H.] probably
\item[I.] sparingly
\item[J.] sprinkle
\item[K.] steer
\item[L.] stimulate
\item[M.] sturdy
\item[N.] synonymous
\item[O.] unique
\end{itemize}

\subsection*{\textbf{Section B}}


\textbf{Directions:} \textit{In this section, you are going to read a passage with ten statements attached to it. Each statement} \textit{contains information given in one of the paragraphs. Identify the paragraph from which the information is derived.} \textit{You may choose a paragraph more than once. Each paragraph is marked with a letter. Answer the questions by} \textit{marking the corresponding letter on }\textbf{\textit{Answer Sheet 2}}\textit{.} \textbf{Treasure Fever}


A) Most visitors come to Cape Canaveral, on the northeast coast of Florida, for the tourist attractions. It’s home to the second-busiest cruise ship port in the world and is a gateway to the cosmos. Nearly 1.5 million visitors flock here every year to watch rockets, spacecraft, and satellites blast off into the solar system from Kennedy Space Center Visitor Complex. Nearly 64 kilometers of undeveloped beach and 648 square kilometers of protected refuge fan out from the cape’s sandy shores.


B) Yet some of Cape Canaveral’s most legendary attractions lie unseen, wedged under the sea’s surface in mud and sand, for this part of the world has a reputation as a deadly ship trap. Over the centuries, dozens of majestic Old World sailing ships smashed and sank on this irregular stretch of windy Florida coast. They were vessels built for war and commerce, crossing the globe carrying everything from coins to cannons, boxes of silver and gold, chests of jewels and porcelain, and pearls from the Caribbean.


C) Cape Canaveral contains one of the greatest concentrations of colonial shipwrecks in the world. In recent years, advances in radar, diving, detection equipment, computers, and GPS have transformed the hunt. The naked eye might see a pile of rocks, but technology can reveal the precious artifacts (人工制品) that lie hidden on the ocean floor.


D) As technology renders the seabed more accessible, the hunt for treasure-filled ships has drawn a fresh tide of salvors (打捞人员) and their investors — as well as marine archaeologists (考古学家) wanting to bring to light the lost relics. But of late, when salvors have found vessels, their rights have been challenged in court. The big question: who should have control of these treasures?


E) High-stakes fights over shipwrecks pit archaeologists against treasure hunters in a vicious cycle of accusations. Archaeologists regard themselves as protectors of history, and they see salvors as careless destroyers. Salvors feel they do the hard work of searching for ships, only to have them stolen from under them when discovered. This kind of clash inevitably takes place on a grand scale. Aside from the salvors, their investors, and the maritime archaeologists who serve as expert witnesses, the battles sweep in local and international governments and organizations like UNESCO that work to protect under-water heritage. The court cases that ensue stretch on for years. Are finders keepers, or do the ships belong to the countries that made them and sent them sailing centuries ago? Where once salvors and archaeologists worked side by side, now they belong to opposing, and equally contemptuous, tribes.


F) Nearly three million vessels lie wrecked on the Earth’s ocean floor — from old canoes to the Titanic — and likely less than one percent have been explored. Some — like an ancient Roman ship found off Antikythera, Greece, dated between 70 and 60 BC and carrying astonishingly sophisticated gears and dials for navigating by the sun — are critical to a new understanding of our past. No wonder there is an eternal stirring among everybody from salvors to scholars to find them.


G) In May 2016, a salvor named Bobby Pritchett, president of Global Marine Exploration (GME) in Tampa, Florida, announced that he had discovered scattered remains of a ship buried a kilometer off Cape Canaveral. Over the prior three years, he and his crew had obtained 14 state permits to survey a nearly 260-square-kilometer area off the cape; they worked 250 days a year, backed by investor funds of, he claims, US \$4 million. It was hard work. Crew members were up at dawn, dragging sensors from their expedition vessels back and forth, day in and day out, year after year, to detect metal of any kind. Using computer technology, Pritchett and his crew created intricate, color-coded maps marked with the GPS coordinates of thousands of finds, all invisible under a meter of sand.


H) One day in 2015, the magnetometer (磁力计) picked up metal that turned out to be an iron cannon; when the divers blew the sand away, they also discovered a more precious bronze cannon with markings indicating French royalty and, not far off, a famous marble column carved with the coat of arms of France, known from historical paintings. The discovery was cause for celebration. The artifacts indicated the divers had likely found the wreck of La Trinity, a 16th-century French vessel that had been at the center of a bloody battle between France and Spain that changed the fate of the United States of America.


I) And then the legal storm began, with GME and Pritchett pitted against Florida and France. The Sunken Military Craft Act of 2004, a US federal act, protects any vessel that was on a military mission, allowing the originating country to claim their ship even centuries later. In 2018, two long years after Pritchett’s discovery, the federal district court ruled in favor of France. For Pritchett, the decision was devastating. Millions of dollars of investor funding and years of labor were lost.


J) But this is far from the first time a salvor has lost all rights to a discovery. In 2012, for instance, Spain won a five-year legal battle against Odyssey Marine Exploration, which had hauled 594,000 gold and silver coins from a Spanish wreck off the coast of Portugal across the Atlantic to the United States. “Treasure hunters can be naive,” says attorney David Concannon, who has had several maritime archaeologists as clients and represented two sides in the battles over the Titanic for 20 years. “Many treasure hunters don’t understand they are going to have to fight for their rights against a government that has an endless supply of money for legal battles that treasure hunters are likely to lose.”


K) Putting an inflated price on artifacts rather than viewing them as cultural and historical treasures that transcend any price is what irritates many archaeologists. For the archaeologist, everything in a wreck matters — hair, fabric, a fragment of a newspaper, rat bones — all things speak volumes. Archaeologists don’t want artifacts ending up in a private collection instead of taking humanity on a journey of understanding.


L) George Bass is one of the pioneers of under-water archaeology, and a researcher at Texas A\&M University. He has testified in court against treasure hunters, but says archaeology is not without its own serious problems. He believes archaeologists need to do a better job themselves instead of routinely criticizing treasure hunters. “Archaeology has a terrible reputation for not publishing enough on its excavations (发掘) and finds,” he says. Gathering data, unearthing and meticulously preserving and examining finds, verifying identity and origin, piecing together the larger story, and writing and publishing a comprehensive paper or book can take decades. A bit cynically, Bass describes colleagues who never published because they waited so long they became ill or died. Who is more at fault, Bass asks, the professional archaeologist who carefully excavates a site and never publishes on it or the treasure hunter who locates a submerged wreck, salvages part, conserves part, and publishes a book on the operation?


M) Pritchett concedes that his find deserves careful excavation and preservation. “I think what I found should go in a museum,” he says. “But I also think I should get paid for what I found.” Indeed, it’s a bit of a mystery why governments, archaeologists, and treasure hunters can’t work together — and why salvors aren’t at least given a substantial finder’s fee before the original owner takes possession of the vessel and its artifacts.


36. Exploration of shipwrecks on the sea floor is crucial in updating our understanding of humanity’s past.


37. Quite a number of majestic ships sailing from Europe to America were wrecked off the Florida coast over


the centuries.


38. Pritchett suffered a heavy loss when a US district court ruled against him.


39. Recently, people who found treasures in shipwrecks have been sued over their rights to own them.


40. Pritchett claims he got support of millions of dollars from investors for his shipwreck exploration.


41. One pioneer marine scientist thinks archaeologists should make greater efforts to publish their findings.


42. With technological advancement in recent years, salvors now can detect the invaluable man-made objects


lying buried under the sea.


43. According to a lawyer, many treasure hunters are susceptible to loss because they are unaware they face a


financially stronger opponent in court.


44. Salvors of treasures in sunken ships and marine archaeologists are now hostile to each other.


45. Archaeologists want to see artifacts help humans understand their past instead of being sold to private


collectors at an outrageous price.


\subsection*{\textbf{Section C}}


\textbf{Directions: }\textit{There are 2 passages in this section. Each passage is followed by some questions or unfinished} \textit{statements. For each of them there are four choices marked A), B), C) and D). You should decide on the best} \textit{choice and mark the corresponding letter on }\textbf{\textit{Answer Sheet 2}}\textit{ with a single line through the centre.}


\subsection*{\textbf{Passage One}}


\textbf{Questions 46 to 50 are based on the following passage.}


Could you get by without using the internet for four and a half years? That’s exactly what singer and actress Selena Gomez has done in a bid to improve her mental health.


She has spoken extensively about the relationship between her social media usage and her mental wellbeing, recalling feeling like “an addict” when she became Instagram’s most followed user in 2016. “Taking a break from social media was the best decision that I’ve ever made for my mental health,” says she. “The unnecessary hate and comparisons went away once I put my phone down.”


Ditching the web at large, however, is a far more subtle and complicated prospect. The increasing digitisation of our society means that everything from paying a gas bill to plotting a route to a friend’s house and even making a phone call is at the mercy of your internet connection. Actively opting out of using the internet becomes a matter of privilege.


Ms Gomez’s multi-millionaire status has allowed her to take the “social” out of social media, so she can continue to leverage her enormous fame while keeping the trolls (恶意挑衅的贴子) at bay . The fact that she’s still the second most-followed woman on Instagram suggests it’s entirely possible to maintain a significant web profile to promote various projects — by way of a dedicated team — without being exposed to the cruel comments, hate mail and rape or death threats.


It goes without saying that this is fundamentally different from how the rest of us without beauty deals and films to publicise use the likes of Instagram, TikTok and Twitter, but even the concept of a digital detox (戒瘾) requires having a device and connectivity to choose to disconnect from.


The UK’s digital divide has worsened over the past two years, leaving poorer families without broadband connections in their homes. Digital exclusion is a major threat to wider societal equality in the UK, so witnessing companies like Facebook championing the metaverse (元宇宙) as the next great frontier when school children are struggling to complete their homework feels particularly irritating.


Consequently, it’s worth bearing in mind that while deleting all social media accounts will undoubtedly make some feel infinitely better, many other people benefit from the strong sense of community that sharing platforms can breed.


Internet access will continue to grow in importance as we edge further towards web 3.0, and greater resources and initiatives are needed to provide the underprivileged with the connectivity they desperately need to learn, work and live. It’s crucial that people who feel that social media is having a detrimental effect on their mental health are allowed to switch off — and for those living in digital exclusion to be able to switch on in the first place.


46. What do we learn about singer and actress Selena Gomez in the past four and a half years?


\begin{itemize}[leftmargin=2em,itemsep=.2em,topsep=.4em]\item[A.] She has had worsening mental problems.
\item[B.] She has won Instagram’s most followers.
\item[C.] She has refrained from using social media.
\item[D.] She has succeeded in a bid on the internet.
\end{itemize}

47. Why does actively opting out of using the internet become a matter of privilege?


\begin{itemize}[leftmargin=2em,itemsep=.2em,topsep=.4em]\item[A.] Most people find it subtle and complicated to give up using the internet.
\item[B.] Most people can hardly ditch the web while avoiding hate and comparisons.
\end{itemize}

\begin{itemize}[leftmargin=2em,itemsep=.2em,topsep=.4em]\item[C.] Most people can hardly get by without the internet due to growing digitisation.
\item[D.] Most people have been seriously addicted to the web without being aware of it.
\end{itemize}

48. Why does the author say “witnessing companies ... feels particularly irritating” (Para. 6)?


\begin{itemize}[leftmargin=2em,itemsep=.2em,topsep=.4em]\item[A.] The UK digital divide would further worsen due to the metaverse.
\item[B.] The concept of the metaverse is believed to be still quite illusory.
\item[C.] School children would be drawn farther away from the real world.
\item[D.] Most families in the UK do not have stable broadband connections.
\end{itemize}

49. What is worth bearing in mind concerning social media platforms?


\begin{itemize}[leftmargin=2em,itemsep=.2em,topsep=.4em]\item[A.] They are conducive to promoting societal equality.
\item[B.] They help many people feel connected with others.
\item[C.] They provide a necessary device for a digital detox.
\item[D.] They create a virtual community on the internet.
\end{itemize}

50. What does the author think is really important for those living in digital exclusion?


\begin{itemize}[leftmargin=2em,itemsep=.2em,topsep=.4em]\item[A.] Having access to the internet.
\item[B.] Edging further towards web 3.0.
\item[C.] Getting more educational resources.
\item[D.] Opening more social media accounts.
\end{itemize}

\subsection*{\textbf{Passage Two}}


\textbf{Questions 51 to 55 are based on the following passage.}


Psychologists have long been in disagreement as to whether competition is a learned or a genetic component of human behavior. Whatever it is, you cannot but recognize the effect competition is exerting in academics and many other areas of contemporary life.


Psychologically speaking, competition has been seen as an inevitable consequence of human drives. According to Sigmund Freud, humans are born screaming for attention and full of organic drives for fulfillment in various areas. Initially, we compete for the attention of our parents. Thereafter, we are at the mercy of a battle between our base impulses for self-fulfillment and social and cultural norms which prohibit pure indulgence.


Current work in anthropology (人类学) has suggested, however, that this view of the role of competition in human behavior may be incorrect. Thomas Hobbes, one of the great philosophers of the seventeenth century, is perhaps best remembered for his characterization of the “natural world”, that is, the world before the imposition of the will of humanity, as being “nasty, brutish, and short.” This image of the pre-rational world is still widely held, reinforced by Charles Darwin’s highly influential work. \textit{The Origin of Species}. which established the doctrine of natural selection. This doctrine, which takes for granted that those species best able to adapt to and master the natural environment in which they live will survive, has suggested that the struggle for survival is an inherent human trait which determines a person’s success. Darwin’s theory has even been summarized as “survival of the fittest” — a phrase Darwin himself never used — further highlighting competition’s role in success. As it has often been pointed out, however, there is nothing in the concept of natural selection that suggests that competition is the most successful strategy for “survival of the fittest.” Darwin asserted in \textit{The Origin of Species} that the struggles he was describing should be viewed as metaphors and could easily include dependence and cooperation.


Many studies have been conducted to test the importance placed on competition as opposed to other values, such as cooperation — by various cultures, and generally conclude that Americans uniquely praise competition as natural, inevitable, and desirable. In 1937, the world-renowned anthropologist Margaret Mead published \textit{Cooperation and Competition among Primitive Peoples}, based on her studies of several societies that did not prize competition, and, in fact, seemed at times to place a negative value on it. One such society was the Zuni Indians of Arizona, and they, Mead found, valued cooperation far more than competition. After studying dozens of such cultures, Mead’s final conclusion was that competitiveness is a culturally created aspect of human behavior, and that its prevalence in a particular society is relative to how that society values it.


51. What does the author think is easy to see in many areas of contemporary life?


\begin{itemize}[leftmargin=2em,itemsep=.2em,topsep=.4em]\item[A.] The disagreement on the inevitability of competition.
\item[B.] The consequence of psychological investigation.
\item[C.] The effect of human drives.
\item[D.] The impact of competition.
\end{itemize}

52. According to psychology, what do people strive to do following the initial stage of their life?


\begin{itemize}[leftmargin=2em,itemsep=.2em,topsep=.4em]\item[A.] Fulfill individual needs without incurring adverse effects of human drives.
\item[B.] Indulge in cultural pursuits while keeping their base impulses at bay.
\item[C.] Gain extensive recognition without exposing pure indulgence.
\item[D.] Satisfy their own desires while observing social conventions.
\end{itemize}

53. What do we learn about the “natural world” characterized by Thomas Hobbes?


\begin{itemize}[leftmargin=2em,itemsep=.2em,topsep=.4em]\item[A.] It gets misrepresented by philosophers and anthropologists.
\item[B.] It gets distorted in Darwin’s \textit{The Origin of Species}.
\item[C.] It is free from the rational intervention of humans.
\item[D.] It is the pre-rational world rarely appreciated nowadays.
\end{itemize}

54. What can we conclude from Darwin’s assertion in \textit{The Origin of Species}?


\begin{itemize}[leftmargin=2em,itemsep=.2em,topsep=.4em]\item[A.] All species inherently depend on others for survival.
\item[B.] Struggles for survival do not exclude mutual support.
\item[C.] Competition weighs as much as cooperation as a survival strategy.
\item[D.] The strongest species proves to be the fittest in natural selection.
\end{itemize}

55. What conclusion did Margaret Mead reach after studying dozens of different cultures?


\begin{itemize}[leftmargin=2em,itemsep=.2em,topsep=.4em]\item[A.] It is characteristic of humans to be competitive.
\item[B.] Americans are uniquely opposed to cooperation.
\item[C.] Competition is relatively more prevalent in Western societies.
\item[D.] People’s attitude towards competition is actually culture-bound.
\end{itemize}

\subsection*{\textbf{Part IV Translation (30 minutes)}}


\textbf{Directions: }\textit{For this part, you are allowed 30 minutes to translate a passage from Chinese into English. You} \textit{should write your answer on }\textbf{\textit{Answer Sheet 2}}\textit{.}


在中国，随着老龄化社会的到来，养老受到普遍关注。人们谈论最多的是应当采取什么样的养老模式。多数人认为，养老模式需要多元化。可以通过政府引导和社会参与，建立更多更好的养老服务机构，改进社区服务中心，鼓励居家自助养老，还可以推行家庭养老与社会养老相结合的模式。随着政府和社会对养老服务事业投入的持续增加，养老设施将不断升级，服务质量逐步改进，老年人的生活将会更加方便舒适、健康快乐。

\end{document}